\documentclass[11pt,a4paper]{article}
\usepackage[margin=1in]{geometry}
\usepackage{amsmath,amssymb,amsfonts}
\usepackage{mathtools}
\usepackage{lmodern}
\usepackage[T1]{fontenc}
\usepackage[utf8]{inputenc}
\usepackage{textcomp}
\usepackage{microtype}
\usepackage{parskip}
\usepackage{enumitem}
\usepackage{array}
\usepackage{booktabs}
\usepackage{longtable}
\usepackage{hyperref}
\hypersetup{hidelinks,
  pdftitle={Paper 087 -- The 13 Primitives: Position Paper of the Event Density Generative Substrate Ontology},
  pdfauthor={Allen Proxmire}}
\setlength{\parindent}{0pt}
\setlength{\parskip}{6pt plus 2pt minus 1pt}

\title{\vspace{-1cm}\Large\bfseries Paper 087 --- The 13 Primitives\\[6pt]
\large\mdseries Position Paper of the Event Density Generative Substrate Ontology}
\author{Allen Proxmire}
\date{May 2026}

\begin{document}
\maketitle

\section*{Preamble: What This Paper Does NOT Claim}

\begin{enumerate}[noitemsep]
\item \textbf{The 13 primitives are not derived.} They are postulated within the Event Density Generative Substrate Ontology. This paper does not attempt to derive them from a deeper structural layer.
\item \textbf{No claim that the substrate is the only possible ontology.} Other substrate ontologies (with different primitive sets) are conceivable; ED commits to this specific 13-primitive set.
\item \textbf{No claim that any downstream result is ``forced from nothing.''} Every downstream result is \emph{conditional} on the 13 primitives + V1/V5 inheritance + upstream papers.
\item \textbf{No claim that the primitives are minimal in some absolute sense.} Whether the 13-primitive set could be reduced to fewer primitives (with the remaining ones derived from a subset) is structurally open. The current commitment is to 13 --- claimed \emph{necessary-for-current-derivations} (removing any one blocks at least one downstream derivation, audited in \texttt{PRIMITIVE\_LOAD\_BEARING\_AUDIT.md}), not minimal in an absolute sense. Two specific minimality caveats: \textbf{P03 bundles three commitments} (channel indexing, locus indexing, and spatial homogeneity) that could in principle be split; and \textbf{the arrow of time is supplied jointly by P11 + P13 + V1's retardation property} (Theorem T18), with the redundancy question among these three structurally open.
\item \textbf{No claim that the primitives capture all of physical reality.} Specific value-layer empirical content (mass spectrum, coupling constants, cosmological parameters) is inherited from empirical measurement, not derived from the primitives.
\item \textbf{No claim that P14 (Bilocal Strain Coupling) is part of the 13.} A P14 \emph{placeholder slot} is reserved (\S4.14) for the bilocal-strain postulate introduced in Paper\_030, but P14 is \textbf{not} included in the canonical 13. Whether P14 should be elevated to the canonical set is open.
\item \textbf{No claim that this paper is a foundational metaphysics.} ED's intellectual neighborhood is the \textbf{substrate-ontology research program} ('t Hooft's cellular-automaton interpretation, Wolfram's Ruliad / hypergraph rewriting, the causal-set program). The methodological precedent for \emph{positing axioms and arguing from downstream reach} runs through Newton, Einstein, Schr\"odinger, Dirac, and the operational-reconstruction tradition (Hardy 2001, CDP 2011, Masanes-M\"uller 2011, Coecke-Kissinger 2017) --- but ED is \textbf{not in the operational-reconstruction lineage}: those programs derive finite-dimensional QM from operational axioms with closed proofs, while ED postulates a substrate and exhibits cross-domain reach without a closed substrate$\to$continuum reconstruction theorem. The primitives are operational substrate-level commitments, evaluated by downstream cross-domain reach.
\item \textbf{No claim that the primitives are testable in isolation.} Individual primitives are not directly empirically testable; the empirical case rests on cross-domain downstream content (substrate-gravity arc, Q-COMPUTE arc, BH arc, etc.).
\end{enumerate}

\section*{Abstract}

The Event Density (ED) substrate is a \textbf{13-primitive generative system}. This paper is the position-paper definition of those 13 primitives --- the canonical enumeration \textbf{P01 through P13} --- together with the postulational-genre framing in which they operate. The primitives are \textbf{postulated, not derived}. ED's closest intellectual neighbors are substrate-ontology research programs --- 't Hooft's cellular-automaton interpretation, Wolfram's Ruliad / hypergraph rewriting, and the causal-set program. The methodological precedent for \emph{positing axioms and arguing from downstream reach} runs through Newton, Einstein, Schr\"odinger, Dirac, and the operational-reconstruction tradition (Hardy 2001, CDP 2011, Masanes-M\"uller 2011, Coecke-Kissinger 2017) --- but ED's epistemic standing is honestly closer to Wolfram's Ruliad than to Hardy's closed-proof reconstruction of finite-dimensional QM. ED postulates a substrate and exhibits cross-domain reach without a closed substrate$\to$continuum reconstruction theorem. All downstream content in the ED corpus --- the gravity arc (Newton's $G$, $a_0$, ECR, BTFR), the QM-kinematics arc (Born rule, Hilbert-space emergence, gauge fields), the BH/Hawking arc, the soft-matter arc, the Q-COMPUTE arc, the entanglement arc --- is \textbf{conditional structural derivation given the 13 primitives + the V1/V5 kernel inheritance}. The empirical case for the substrate is the \textbf{cross-domain reach} of these conditional derivations across regimes that share no obvious common physics in standard frameworks. The paper makes no claim that the primitives are forced from a deeper structural layer, no claim that ED is the only consistent substrate ontology, no claim that the primitive set is minimal in some absolute sense, and no claim of derivation from nothing. A P14 \emph{placeholder slot} is reserved for the Bilocal Strain Coupling postulate introduced in Paper\_030; P14 is \textbf{not} included in the canonical 13 at this time.

\section{Introduction}

\subsection{What this paper does}

This paper establishes the \textbf{authoritative Wave-2 definition} of the 13-Primitive Generative System. It serves three functions:

\begin{enumerate}[noitemsep]
\item \textbf{Canonical enumeration.} Each of P01 through P13 is stated, named, and given its operational content.
\item \textbf{Postulational-genre framing.} The 13 primitives are postulated; downstream content is conditional structural derivation. No forcing-era language.
\item \textbf{Cross-paper anchor.} Every other Wave-2 generative paper references the position paper for primitive enumeration; this paper is the single source of truth.
\end{enumerate}

It is the \textbf{position-paper anchor} for the ED Generative Primitives System, written after the pivot from the M-series ``primitive forcing'' program (which external review identified as circular).

\subsection{What this paper is NOT}

This paper is \emph{not}:

\begin{itemize}[noitemsep]
\item A derivation of the primitives from a deeper layer.
\item An empirical defense of the primitives (the empirical case is cross-domain reach, developed in downstream papers).
\item A historical narrative of how ED arrived at this primitive set.
\item A philosophical metaphysics.
\item A complete operational manual (specific operational protocols for each primitive are developed in subsequent papers).
\end{itemize}

\subsection{The genre: conditional structural derivation}

The genre of Wave-2 ED is \textbf{conditional structural derivation given postulated primitives}. The structure is:

\begin{quote}
\emph{Given the 13 primitives (+ V1/V5 inheritance + upstream paper results), the following structural consequence follows: [downstream result].}
\end{quote}

\textbf{Intellectual neighborhood (substrate-ontology research programs):}

\begin{itemize}[noitemsep]
\item \textbf{'t Hooft cellular-automaton interpretation of QM:} discrete substrate posited; quantum mechanics emerges as a coarse-grained description.
\item \textbf{Wolfram Ruliad / hypergraph-rewriting program:} computational substrate posited; physics emerges as exhibits of substrate-level structure.
\item \textbf{Causal-set program (Sorkin et al.):} discrete causal-set substrate posited; continuum spacetime emerges via coarse-graining.
\end{itemize}

ED sits structurally adjacent to these programs: ontology + cross-domain exhibits, not an axiomatic reconstruction.

\textbf{Methodological precedent (positing axioms and arguing from downstream reach):}

\begin{itemize}[noitemsep]
\item \textbf{Newton's three laws:} posited; classical mechanics follows.
\item \textbf{Einstein's two postulates of special relativity:} posited; relativistic kinematics follows.
\item \textbf{Schr\"odinger's equation:} posited; non-relativistic QM follows.
\item \textbf{Dirac equation:} posited under Lorentz-covariance + first-order requirements; spinor structure and $g = 2$ follow.
\end{itemize}

\textbf{A different genre --- operational reconstruction:}

\begin{itemize}[noitemsep]
\item \textbf{Hardy 2001 / Masanes-M\"uller 2011 / CDP 2011:} operational postulates posited; finite-dimensional QM reconstructed by \textbf{closed proofs}.
\item \textbf{Coecke--Kissinger 2017:} symmetric-monoidal-dagger-category axioms posited; categorical QM follows by closed-proof reconstruction.
\end{itemize}

\textbf{ED is not in the operational-reconstruction lineage.} Those programs derive finite-dimensional QM from operational axioms by closed proofs; ED postulates a substrate and exhibits cross-domain reach without a closed substrate$\to$continuum reconstruction theorem. The honest comparison places ED in the Wolfram-Ruliad / 't Hooft / causal-set lineage. The methodological precedent (Newton, Einstein, etc.) is invoked as a posit-then-argue-downstream genre, not as a claim of equivalent epistemic standing with closed-proof reconstructions.

The empirical case rests on cross-domain reach (what the postulates \emph{generate} across domains), not on derivation of the postulates themselves.

\section{Primitive Inputs}

\emph{This paper's ``primitive inputs'' are themselves the canonical primitives --- there is no upstream substrate layer. The 13 primitives are postulated within the ED Generative Substrate Ontology.}

\textbf{Upstream dependencies for the position-paper genre framing:}

\begin{itemize}[noitemsep]
\item Substrate-ontology research-program lineage ('t Hooft cellular-automaton interpretation, Wolfram Ruliad / hypergraph rewriting, causal-set program). The closest neighbors to ED.
\item Methodological precedent (Newton, Einstein, Schr\"odinger, Dirac) for \emph{positing axioms and arguing from downstream reach}. Not an equivalent epistemic-standing claim; ED is closer to Wolfram-Ruliad than to closed-proof reconstructions.
\item Distinguished from (not in) the operational-reconstruction tradition (Hardy, CDP, Masanes-M\"uller, Coecke-Kissinger) --- see \S1.3.
\item The M-series archive which documents the pre-pivot exploratory ``primitive forcing'' attempt that external review identified as circular.
\end{itemize}

\textbf{No empirical inputs at the primitive level.} The empirical case is downstream cross-domain reach.

\section*{\S2.5 Load-Bearing Step Audit}

\begin{center}
\footnotesize
\begin{tabular}{p{0.34\textwidth}p{0.14\textwidth}p{0.46\textwidth}}
\toprule
\textbf{Step} & \textbf{Status} & \textbf{Source / justification} \\
\midrule
The 13 primitives are postulated & P & This paper's foundational commitment \\
Substrate-ontology research-program lineage ('t Hooft, Wolfram, causal-set) + postulational-genre methodological precedent (Newton, Einstein, Schr\"odinger, Dirac) & I & Inherited from substrate-ontology and postulational traditions; \textbf{distinguished from operational-reconstruction tradition} (Hardy, CDP, etc.) --- see \S1.3 \\
Each $P_i$'s operational content & P & \S4 below; each is a substrate-level commitment \\
Downstream results are conditional & P (genre tautology) & Follows tautologically from the postulated conditional-derivation genre; not derived from primitives via algebraic/logical operations. \\
Cross-domain reach is the empirical case & A$\to$P & \textbf{Position commitment} of this paper; not derived. Moved to \S6.4 explicitly. \\
P14 placeholder reserved but not canonical & P & \S4.14 explicit reservation \\
No primitive forcing & P (genre tautology) & Tautological consequence of the postulational genre commitment. \\
\bottomrule
\end{tabular}
\end{center}

All load-bearing steps are P (postulational position commitments) or D (structural consequences of the postulational genre) or I (inherited from tradition). \textbf{No A (asserted) rows.} Paper is publishable.

\section{The Postulational Tradition}

\subsection{What ``postulated'' means in ED}

A substrate primitive is \textbf{postulated} if and only if:

\begin{enumerate}[noitemsep]
\item It is stated as a substrate-level commitment without derivation from a deeper layer.
\item Its operational content is specified explicitly (what the primitive \emph{does} at the substrate level).
\item Downstream content is \emph{conditional} on the primitive: removing it changes or breaks downstream derivations.
\item It is evaluated empirically through downstream cross-domain reach, not through direct test of the primitive in isolation.
\end{enumerate}

\subsection{Why postulational, not foundational}

ED does not claim to derive the substrate from deeper physics, philosophy, or logic. The substrate is the \textbf{operational floor} of the framework --- the level at which structural argument begins. Asking ``why these primitives and not others?'' is analogous to asking ``why these axioms of arithmetic and not others?'' within Peano arithmetic: the question is meta-axiomatic, addressable in a different framework, but not addressable within ED.

This is structurally distinct from the M-series pre-pivot program, which attempted to \emph{force} the primitives from upstream M-postulates. External review identified that program as circular: M-1 (the foundational forcing argument) presupposed structures that Paper \#1 of the conditional-derivation corpus then ``derived.'' The Wave-2 pivot dissolves the circularity by declaring the primitives postulated and the downstream content conditional.

\subsection{The empirical case is downstream}

The empirical case for the ED 13-primitive system is:

\begin{itemize}[noitemsep]
\item \textbf{Cross-domain reach:} the same 13 primitives (+ V1/V5 kernel inheritance) generate substrate-level mechanisms for gravity (Newton, $a_0$, ECR, BTFR), QM kinematics (Born rule, Hilbert space, gauge), BH physics (horizon decoupling, Hawking, remnant), soft-matter (Maxwell, DCGT), Q-COMPUTE (multiplicity wall), entanglement (Schmidt, monogamy, no-signaling). The cross-domain coherence is the empirical wedge.
\item \textbf{Falsification via downstream content:} if any downstream conditional derivation is empirically falsified at a level that traces to a specific primitive, that primitive is challenged. The primitives are not directly testable; they are tested through their downstream content.
\end{itemize}

\section{The 13 Primitives --- Canonical Enumeration}

Each primitive is stated as a substrate-level operational commitment. Where downstream papers depend load-bearingly on a primitive, the cross-reference is noted.

\subsection*{P01 --- Event-density layer existence}

A pre-quantum substrate exists as a primitive structural layer. The substrate is not a Hilbert space, not a smooth manifold, not a field theory --- it is a discrete graph-like structure equipped with primitive scalar, relational, and angular quantities. P01 is the foundational \textbf{ontological-existence commitment} that grounds every paper in the program.

The substrate is not derived from a deeper structural layer; the event-density layer is \emph{what there is} at substrate level. Within this layer, the further primitive structures (P02--P13) operate --- including chains as derived composite structures (persistent forward-causal sequences of substrate events within the event-density layer; introduced in P02 as the things that participate).

\emph{Operational content:} a substrate layer exists, with primitive structural content articulated by P02--P13. Below this layer there is no further substrate structure resolved by ED; P01 is the foundational existence commitment.

\emph{Load-bearing in:} every Wave-2 paper. P01 is \textbf{universally implicit} --- without P01, there is no substrate, and the rest of the ontology has no carrier.

\subsection*{P02 --- Participation as primitive relation}

A \textbf{chain} $C$ --- a persistent forward-causal sequence of substrate events within the event-density layer of P01 --- \textbf{participates} in a channel $K$ at substrate locus $u$ at substrate time $t$. The participation relation is a primitive: not derived from a deeper interaction structure, not derived from field-coupling, not derived from gauge connection. Participation is the substrate's primitive mechanism for chain-channel coupling. Chains themselves are \emph{not} canonical primitives separate from P01; they are the substrate-level entities (within the event-density layer) whose participation in channels P02 declares as primitively relational.

\emph{Operational content:} $(C, K, u, t)$ --- the four-tuple parameterizing a substrate participation event. Chains are the derived composite structures (persistent micro-event sequences) that participate; P02 declares the participation relation between chains and channels as primitive.

\emph{Load-bearing in:} all kernel-related papers (089, 090, 093), gravity arc (025, 027), Q-COMPUTE arc (054, 056). Every Wave-2 paper using the language of ``chains,'' ``channels,'' ``participation measure,'' ``participation manifold'' rests on P02.

\subsection*{P03 --- Channel and locus indexing; spatial homogeneity}

The substrate carries discrete index sets:

\begin{itemize}[noitemsep]
\item A \textbf{channel index set} $\mathcal{K}$ over which channels are enumerated.
\item A \textbf{locus index set} indexed by substrate-graph positions $u$.
\end{itemize}

The participation graph (formed by chains, channels, loci) is \textbf{spatially homogeneous}: substrate-level operations are invariant under translations of the locus index. No locus is privileged.

\emph{Operational content:} discrete indexing + translation-invariance. Generates conservation laws via Noether-analog substrate arguments in downstream papers.

\emph{Load-bearing in:} 025 (holographic counting), 027 (Newton's $G$), 029 (cosmic decoupling).

\subsection*{P04 --- Bandwidth as non-negative additive scalar}

Each channel-locus participation carries a \textbf{bandwidth} $b_K(u) \geq 0$. Bandwidth is:

\begin{itemize}[noitemsep]
\item Non-negative.
\item Additive under channel decomposition: $b_{K_1 \cup K_2} = b_{K_1} + b_{K_2}$ for disjoint sub-channels.
\item The substrate-level scalar carrier of ``amount of participation.''
\end{itemize}

\emph{Operational content:} $b_K(u, t) \in \mathbb{R}_{\geq 0}$.

\emph{Load-bearing in:} 089 (V1 kernel envelope), 090 (V5 cross-chain), 054 (UR-1), 056 (Class-A wall), all gravity papers.

\subsection*{P05 --- Polarity-transport along edges}

The substrate's participation graph carries a \textbf{connection structure} that transports \textbf{polarity} along graph edges. Polarity-transport is the substrate-level mechanism for gauge content --- a forward-causal connection-like operation on the participation graph.

\emph{Operational content:} $\pi_K(u, t) \in U(1)$ transports along edges $(u, t) \to (u', t')$ with a substrate-level connection.

\emph{Load-bearing in:} 089 (V1 gauge-compatibility), 090 (V5 gauge-compatibility), QFT arc (T17 gauge fields, paper \#5 / future paper\_015).

\subsection*{P06 --- Spatial dimension primitive ($D = 3+1$)}

The substrate has \textbf{spatial dimension three} (plus one temporal direction via P13). The dimensionality is a substrate primitive --- not derived from compactification of higher-dimensional structure, not derived from a counting argument, not derived from anthropic reasoning.

\emph{Operational content:} substrate-graph adjacency in 3 spatial directions; closed 2-surfaces have area $4\pi R^2$; etc.

\emph{Load-bearing in:} 025 (holographic count area scaling), 027 (Newton inverse-square via $R^2$ surface), NS arc ($D=3+1$ PDE forcing).

\textbf{Acknowledgment --- P06 is doing extraordinary load-bearing work across the corpus.} Among the 13 primitives, P06 is among the most consequential:

\begin{itemize}[noitemsep]
\item Paper\_025's holographic bound's $4\pi R^2$ area scaling depends on P06.
\item Paper\_027's Newton inverse-square law's $R^2$ surface area comes from P06.
\item Paper\_025's codimension-1 channel-footprint argument (P-Codim-1 postulate) is dimension-dependent and operates within P06.
\item All gravity-arc spatial-geometric content (Papers\_027/028/029/030/031/039) implicitly depends on P06.
\end{itemize}

\textbf{Other operational reconstruction programs (Hardy 2001, CDP 2011, Masanes--M\"uller, Coecke--Kissinger) treat spatial dimension as a derived structural fact}, with sometimes-elaborate arguments for why it's 3 (information-theoretic constraints, Lie-algebra-rank arguments, observability-based reconstructions). ED postulates it --- the postulational-tradition move (Newton, Einstein). This is defensible, but the load-bearingness deserves explicit recognition: \textbf{whether P06 is derivable from a subset of the other 12 primitives is structurally open}. If it is, the canonical set shrinks; if not, P06 stands as one of the most consequential commitments of the ED framework, and the corpus's empirical case rests heavily on P06 holding.

\subsection*{P07 --- Channel structure as ontological primitive}

Channels are structurally distinguishable carriers with \textbf{intrinsic identities} in the participation graph. Two distinct channels at the same locus are substrate-level distinct objects, even if their bandwidth and polarity content happen to coincide. The channel structure is an ontological primitive.

\emph{Operational content:} $K \in \mathcal{K}$, with $K_1 \neq K_2$ even if $b_{K_1}(u) = b_{K_2}(u)$ and $\pi_{K_1}(u) = \pi_{K_2}(u)$.

\emph{Load-bearing in:} 025 (channel-counting), 054 (UR-1 multi-channel structure), 089/090 (kernel rule-types).

\subsection*{P08 --- Substrate scale $\ell_\mathrm{ED}$}

The substrate has a characteristic \textbf{edge length / discretization scale} $\ell_\mathrm{ED}$. This scale is a primitive --- the substrate is not continuous below $\ell_\mathrm{ED}$, and substrate channels are resolved at this granularity.

\emph{Operational content:} the smallest substrate-graph separation; below $\ell_\mathrm{ED}$ the substrate has no further resolved structure.

\emph{Empirical identification (inherited, not derived):} $\ell_\mathrm{ED} = \ell_P$ (Planck length) via Newton-recovery in Paper\_027 \S5.3. This identification is \textbf{value-layer empirical content}, not a substrate-level derivation.

\emph{Load-bearing in:} 025, 027, 029, 039, 047, 089, 090.

\subsection*{P09 --- $U(1)$-valued polarity}

Each channel-locus participation carries a \textbf{polarity} $\pi_K(u) \in U(1) \cong S^1$. Polarity is the substrate's primitive angular variable --- single $U(1)$ at the most fundamental level (higher-rank gauge structure emerges via T17 rule-type composition, not via larger-$U(N)$ polarity).

\emph{Operational content:} $\pi_K(u, t) \in [0, 2\pi)$ (or equivalently $e^{i\pi_K} \in U(1)$).

\emph{Load-bearing in:} 089/090 (gauge compatibility), QFT arc, entanglement arc.

\subsection*{P10 --- Rule-type primitive}

The substrate supports \textbf{multiple structurally distinct rule-types}: matter rule-types (chain-carrying), gauge rule-types (T17 bundles), kernel rule-types (V1, V5, future cascade). Rule-type multiplicity is a substrate primitive --- not derived from a deeper unification, not derived from a single ``master rule.''

\emph{Operational content:} $\tau_\bullet$ --- each rule-type a primitive label with its own participation measure / kernel content.

\emph{Load-bearing in:} 089 (V1 rule-type), 090 (V5 rule-type), QFT arc (gauge rule-types).

\subsection*{P11 --- Commitment irreversibility}

At certain substrate-level events (``commitment events''), a chain's multi-channel participation collapses to single-channel participation, with the un-selected channels' phase content randomized. \textbf{The collapse is irreversible}: no substrate-level operation maps post-commitment to pre-commitment state.

\emph{Operational content:} discrete, environmentally-phase-randomizing, irreversible substrate events.

\emph{Load-bearing in:} 089/093 (kernel-arrow T18), 054 (UR-1.3 commitment-injection), 039 (BH information not crossing horizon), all kinematic-arrow content.

\subsection*{P12 --- Stability landscape}

Each chain carries a substrate-level functional:
\[
\Sigma_C = \mathrm{Coh}(C) - \mathrm{Str}(C) - \mathrm{Grad}(C)
\]
where $\mathrm{Coh}$, $\mathrm{Str}$, $\mathrm{Grad}$ are the chain's coherence, strain, and gradient content. The chain's experienced acceleration is the negative gradient of $\Sigma_C$ along its adjacency direction.

\emph{Operational content:} $\Sigma_C: \mathrm{config\text{-}space}(C) \to \mathbb{R}$, with $\vec{a}_C = -\vec{\nabla}_\mathrm{adj} \Sigma_C$.

\emph{Load-bearing in:} 027 (Newton's $G$ via cumulative strain), 029 ($a_0$), 030 (ECR), 039 (BH stability-landscape gradient at horizon).

\subsection*{P13 --- Time homogeneity}

The substrate's dynamical primitive is \textbf{invariant under time translation}: the substrate-level laws at time $t$ are identical to those at time $t'$. No moment is privileged.

\emph{Operational content:} $\partial_t$ is a substrate-level symmetry. Generates energy-conservation-analog substrate content.

\emph{Load-bearing in:} all dynamical papers; ensures P11's commitment-direction is the only temporal-asymmetry primitive.

\subsection{4.14 --- P14 placeholder slot (reserved, not defined here)}

A \textbf{P14 placeholder slot} is reserved for a future substrate-level postulate. The leading candidate is:

\begin{quote}
\textbf{P14 (Bilocal Strain Coupling).} In the joint weak-gradient regime, bilocal substrate channels carry strain proportional to the geometric mean of the local and horizon-sourced bandwidth contributions. \emph{(Articulated in Paper\_030 \S2; not included in the canonical 13.)}
\end{quote}

P14 is \textbf{not defined in this paper} and is \textbf{not included in the canonical 13}. Whether P14 should be elevated to the canonical set, or whether it can be derived from the existing 13 (potentially with additional substrate-level analysis), is structurally open. See \S6 below.

\subsection{4.15 --- Paper-Specific Postulate Inventory}

The Wave-2 corpus introduces \textbf{paper-specific postulates} beyond the canonical 13 as needed for individual derivations. This subsection inventories them in one place.

\textbf{Each postulate is labeled:}
\begin{itemize}[noitemsep]
\item Which paper introduced it.
\item Its role: \textbf{(S) structural} / \textbf{(R) regime} / \textbf{(D) definitional} / \textbf{(C) choice between substrate-consistent models}.
\item Whether it is expected to remain \textbf{(perm) permanent} in the corpus, or \textbf{(prov) provisional} (may be replaced by future derivation from canonical 13 or refinement).
\end{itemize}

\textbf{Inventory:}

{\footnotesize
\begin{longtable}{p{0.27\textwidth}p{0.13\textwidth}p{0.40\textwidth}p{0.12\textwidth}}
\toprule
\textbf{Postulate} & \textbf{Paper} & \textbf{Role} & \textbf{Status} \\
\midrule
\endhead
P14 (Bilocal Strain Coupling) & 030 \S2 & S --- bilocal substrate-channel strain coupling in joint weak-gradient regime & (prov) \\
P-Codim-1 (Channel codimension-1 structure) & 025 \S2 & S --- channels as codimension-1 substrate-graph objects & (perm) \\
P-Sat (Substrate-saturation regime) & 025 \S2 & R --- substrate-saturated regime for bound-equality applications & (perm) \\
P-NoBackwardChain (Backward-chain non-constructibility) & 089 \S2 & S --- backward chain content not constructible from canonical 13 primitives & (perm) \\
P-SubstrateLocality (Substrate-graph locality) & 089 \S2 & S --- substrate operations have bounded range at $\ell_\mathrm{ED}$ scale & (prov) \\
P-Potential-Reading (Cumulative-strain potential reading) & 026 \S2 & C --- choice between Model A (potential) and Model B (per-channel modulation) & (perm) \\
P-V5-Even (V5 envelope parity-symmetric in $\omega$) & 047 \S2 & S --- V5 envelope parity selection enabling $(\omega/\omega_c)^2$ leading correction & (perm) \\
P-LinRate (Substrate commit-rate linear in bandwidth) & 003 \S2 & S --- substrate-level commitment-rate linear in $b_K$ & (prov) \\
P-Rigidity-Gap (Substrate energy gap from global rigidity) & 057 \S2 & S --- global geometric rigidity creates substrate P12 stability-landscape gap & (prov) \\
P-Boltz-Analogy (Substrate Boltzmann-like suppression in gapped regimes) & 057 \S2 & S --- substrate-level transition rates exponentially suppressed by gap & (prov) \\
P-Corr-Budget (Substrate cross-chain correlation budget finite) & 058 \S2 & S --- V5 cross-chain correlation content bounded above & (perm) \\
P-Redundancy-Mapping (Substrate redundancy $\to$ linear coverage) & 058 \S2 & R --- linear $R_C \to$ coverage mapping in unsaturated regime & (perm) \\
P-V5-EntBudget (V5 entanglement-bandwidth budget finite) & 050 \S2 & S --- V5 cross-chain entanglement-amplitude bounded above & (perm) \\
P-Re-routing (Substrate-level entanglement re-routing on saturation) & 050 \S2 & S --- post-Page-time substrate-level re-routing through radiation channels & (perm) \\
P-Fiber-Naming (Channels as fiber identifications) & 015 \S2 & D --- substrate-ontological naming: channels play fiber role in substrate fiber-bundle analog & (perm) \\
P-Connection-Naming (Polarity-transport as substrate connection) & 015 \S2 & D --- substrate-ontological naming: polarity-transport plays connection role & (perm) \\
P-StructureGroup-Naming ($U(1)$ as substrate structure group) & 015 \S2 & D --- substrate-ontological naming: $U(1)$ as structure group at primitive level & (perm) \\
P-Bundle-Definition (Substrate rule-type bundle defined by parts) & 015 \S2 & D --- substrate-level definitional commitment for the bundle object & (perm) \\
P-NonAbelian-Analogy & 015 \S2 & S --- substrate-level analogy to standard non-abelian gauge structure & (prov) \\
P-Polarity-Connection-Match & 016 \S2 & S --- substrate-level identification of DCGT-coarse-grained substrate polarity-transport with continuum gauge potential & (perm) \\
P-Motif-Algebra (Complex vector space on amplitudes) & 007 \S2 & D --- substrate-level definitional commitment: amplitude content forms complex vector space & (perm) \\
P-Bipartite-Mapping (Bipartite-chain substrate-level mapping) & 063 \S2 & S --- substrate-level mapping for two-chain joint state & (perm) \\
\bottomrule
\end{longtable}
}

\textbf{Inventory summary:} 22 paper-specific postulates introduced beyond canonical 13. Of these:

\begin{itemize}[noitemsep]
\item \textbf{9 are (perm)} structural commitments expected to remain permanent in the corpus.
\item \textbf{8 are (perm)} definitional / regime / naming commitments expected to remain permanent.
\item \textbf{5+ are (prov)} provisional postulates that may be replaced by future derivation from canonical 13 or refinement (P14, P-SubstrateLocality, P-LinRate, P-Rigidity-Gap, P-Boltz-Analogy, P-NonAbelian-Analogy).
\end{itemize}

\textbf{The inventory is updated as new Wave-2 papers introduce new postulates.}

\textbf{Open structural question:} the canonical 13 + $\sim$22 paper-specific postulates is a \textbf{postulate set of $\sim$35 commitments}. Whether the substrate-level program can reduce postulate count (via derivability analyses) is the structural-refinement frontier. Currently the count grows as papers are written; whether it can decrease is the durability test of the framework.

\section{The Conditional-Derivation Genre}

\subsection{Structure of a conditional derivation}

A Wave-2 generative paper has the form:

\begin{quote}
\emph{Given the 13 primitives (+ V1/V5 inheritance + upstream paper results + regime conditions), the structural consequence is [result].}
\end{quote}

Each downstream result identifies its load-bearing primitives. Removing or modifying a load-bearing primitive changes or breaks the derivation.

\subsection{What ``given'' means}

``Given'' denotes \textbf{postulational conditioning}, not logical entailment from a deeper layer. The primitives are taken as substrate-level facts; the derivation proceeds from there.

\subsection{What ``structural consequence'' means}

A structural consequence is a downstream result whose form follows from the substrate primitives via:

\begin{itemize}[noitemsep]
\item Substrate-level operational sequences (chains, channels, polarity-transport, commitment).
\item Kernel-mediated content (V1, V5, future cascade).
\item Coarse-graining via DCGT (when hydrodynamic-window scale separation holds).
\item Standard mathematical operations (linear algebra, calculus, dimensional analysis).
\end{itemize}

Structural consequences are derived; their \emph{form} is fixed by the primitives. Their \emph{numerical values} may be inherited from value-layer empirical content (mass spectrum, $\hbar$, $c$, $H_0$, etc.).

\subsection{Cross-domain reach as the empirical case}

The empirical case for the 13 primitives is \textbf{cross-domain reach} --- explicitly \textbf{position-committed in this paper}, not derived:

The same 13 primitives generate substrate-level mechanisms for:

\begin{itemize}[noitemsep]
\item \textbf{Gravity:} Newton's $G$, transition acceleration $a_0$, ED Combination Rule, BTFR slope-4.
\item \textbf{QM kinematics:} Hilbert-space emergence, Born rule, gauge fields (T17), Berry/Bloch/AB phase.
\item \textbf{BH physics:} horizon decoupling, Hawking spectrum, Planck-mass remnant, information paradox not generated.
\item \textbf{Soft matter:} Maxwell viscoelastic relaxation, NS-Smoothness, MHD.
\item \textbf{Q-COMPUTE:} Class-A/B/C multiplicity walls, UR-1.
\item \textbf{Entanglement:} Schmidt, monogamy, no-signaling, von Neumann entropy.
\end{itemize}

The cross-domain coherence --- one substrate primitive set generating mechanisms across these regimes --- is the \textbf{empirical wedge}. This is a position commitment of the ED program, evaluated by external review of the downstream content, not derivable from a deeper layer.

\section{The P14 Placeholder Slot}

\subsection{Why a placeholder, not a canonical primitive}

Paper\_030 (ED Combination Rule) introduced P14 (Bilocal Strain Coupling) as a substrate-level postulate. Whether P14 belongs in the canonical 13 is structurally open:

\begin{itemize}[noitemsep]
\item \textbf{Argument for inclusion:} P14 is load-bearing in Paper\_030 + Paper\_031 (BTFR). If load-bearing primitives belong in the canonical set, P14 should be included.
\item \textbf{Argument against inclusion:} P14 may be derivable from the existing 13 + additional substrate-level analysis (e.g., from the bilocal-channel structure that P02 + P07 + P03 imply). Until that derivation is attempted, P14's status is uncertain.
\end{itemize}

\subsection{Operational stance}

The current operational stance is:

\begin{itemize}[noitemsep]
\item \textbf{The canonical 13} (P01--P13) defined in \S4 is the baseline.
\item \textbf{P14 placeholder slot} is reserved (this paper \S4.14) for the Bilocal Strain Coupling postulate.
\item \textbf{Papers using P14} (Paper\_030, Paper\_031) cite P14 with explicit conditional dependency.
\item \textbf{Future work} may either (a) derive P14 from the canonical 13, demoting the placeholder to a theorem, or (b) elevate P14 to the canonical set (becoming the 14-Primitive Generative System).
\end{itemize}

The placeholder is open until resolved.

\subsection{Other placeholder slots}

No other placeholder slots are reserved at this time. If future Wave-2 work identifies additional load-bearing substrate-level postulates that are not derivable from the canonical 13, new placeholder slots may be added (P15, P16, etc.), each labeled openly.

\section{Cross-Paper Anchor Role}

Every Wave-2 generative paper references this position paper for the canonical primitive enumeration. The convention is:

\begin{quote}
\emph{``This paper takes $P_n$, $P_m$, $\ldots$ as postulated within the ED Generative Primitives System (see Paper\_087).''}
\end{quote}

This paper is the \textbf{single source of truth} for the primitive list. Any future revisions to the canonical 13 (or to the placeholder slot P14) propagate from edits to this paper.

\textbf{Cross-reference inventory (existing Wave-2 papers):}

\begin{itemize}[noitemsep]
\item Paper\_027 (Newton's $G$): uses P02, P03, P04, P07, P08, P12, P13 + V1.
\item Paper\_029 (transition acceleration $a_0$): uses P02--P04, P07, P08, P11--P13 + V1.
\item Paper\_030 (ED Combination Rule): uses P02--P04, P07, P08, P11--P13 + V1 + \textbf{P14 placeholder}.
\item Paper\_031 (BTFR slope-4): inherits from 027, 029, 030 (and through 030, P14 placeholder dependency).
\item Paper\_039 (Horizon Decoupling): uses P02, P04, P05, P07, P09, P10, P11, P12 + V1 + V5.
\item Paper\_056 (Class-A Wall): uses P02, P04, P05, P07, P09, P10, P11, P12 + V1 + UR-1 (Paper\_054, in queue).
\item Paper\_090 (V5 Cross-Chain Kernel): uses P02, P04, P05, P07, P09, P10, P11 + V1.
\item Paper\_093 (Kernel Arrow T18): uses P02, P04, P05, P07, P09, P10, P11, P13 + V1 admissible class.
\end{itemize}

No paper currently uses all 13; load-bearing subsets vary by domain.

\section{Falsification Criteria}

The 13 primitives are not directly testable. Falsification is via downstream content. The following observations would falsify the 13-primitive system:

\subsection{Downstream conditional-derivation failure traced to a specific primitive}

If a Wave-2 conditional derivation is empirically falsified, and the failure traces to a specific primitive (e.g., empirical evidence of advanced kernel propagation traces to P11), the load-bearing primitive is challenged.

\textbf{Falsifier sentence:} \emph{Empirical demonstration that any substrate-level forward-causal commitment-irreversibility prediction (decoherence directionality, BH information not escaping, kernel-level arrow) fails systematically would falsify P11.}

\subsection{Cross-domain reach failure}

If the 13 primitives generate empirically-correct content in some regimes but systematically empirically-incorrect content in others, the cross-domain reach claim is falsified, and the substrate is structurally challenged.

\textbf{Falsifier sentence:} \emph{Demonstration that the substrate framework produces predictions consistent with empirical content in gravity/QM/BH but systematically inconsistent in soft-matter/Q-COMPUTE/entanglement (or any reverse pattern) would falsify the cross-domain reach as the empirical case.}

\subsection{Derivability of the primitives from a deeper layer}

If a deeper structural framework is identified that \emph{derives} the 13 primitives (rather than postulating them), the postulational genre is superseded. ED would then become a \emph{theorem} of the deeper framework rather than a postulational substrate ontology.

\textbf{Falsifier sentence:} \emph{Identification of a deeper structural framework (mathematical, physical, or operational) that derives P01--P13 from a more fundamental layer would supersede the postulational stance, repositioning ED as a theorem rather than an axiom set.}

This is not falsification of the \emph{content} of the primitives --- they would presumably remain operationally correct --- but of the \emph{positional} status of ED as the foundational level.

\subsection{Primitive minimality violation}

If a strict subset of the 13 primitives is shown to derive the remaining ones, the primitive set is non-minimal and should be reduced.

\textbf{Falsifier sentence:} \emph{Demonstration that the canonical 13 contains a derivable subset (e.g., P13 derivable from P11 + P03, or P06 derivable from P03 + P07) would reduce the primitive count without changing downstream content.}

\subsection{P14 status resolution}

If P14 (Bilocal Strain Coupling) is derived from the canonical 13, the placeholder slot is eliminated and Paper\_030's framing is updated. If P14 is shown to be non-derivable and load-bearing across additional domains, P14 is elevated to the canonical set (becoming P14, 14-Primitive Generative System).

\textbf{Falsifier sentence:} \emph{Empirical evidence that no bilocal-strain-coupling content of any form is required by the substrate (i.e., the joint weak-gradient regime can be handled by the canonical 13 alone) would resolve the P14 placeholder without elevation.}

\section{Position Statement}

The Event Density (ED) substrate is a \textbf{13-Primitive Generative System} situated in the substrate-ontology research-program lineage ('t Hooft cellular-automaton interpretation, Wolfram Ruliad / hypergraph rewriting, causal-set program). The methodological precedent for \emph{positing axioms and arguing from downstream reach} runs through Newton, Einstein, Schr\"odinger, Dirac, and the operational-reconstruction tradition (Hardy 2001, CDP 2011, Masanes-M\"uller 2011, Coecke-Kissinger 2017) --- but ED's epistemic standing is honestly closer to Wolfram's Ruliad than to closed-proof operational reconstructions; ED postulates a substrate and exhibits cross-domain reach without a closed substrate$\to$continuum reconstruction theorem. The canonical 13 primitives are:

\begin{quote}
\textbf{P01} event-density layer existence;\\
\textbf{P02} participation as primitive relation (chains participating in channels);\\
\textbf{P03} channel + locus indexing with spatial homogeneity;\\
\textbf{P04} bandwidth as non-negative additive scalar;\\
\textbf{P05} polarity-transport along edges;\\
\textbf{P06} spatial dimension $D = 3+1$;\\
\textbf{P07} channel structure as ontological primitive;\\
\textbf{P08} substrate scale $\ell_\mathrm{ED}$;\\
\textbf{P09} $U(1)$-valued polarity;\\
\textbf{P10} rule-type primitive;\\
\textbf{P11} commitment irreversibility;\\
\textbf{P12} stability landscape;\\
\textbf{P13} time homogeneity.
\end{quote}

A \textbf{P14 placeholder slot} is reserved for the Bilocal Strain Coupling postulate articulated in Paper\_030; P14 is \textbf{not} in the canonical 13 at this time.

All downstream content in the ED corpus is \textbf{conditional structural derivation given the 13 primitives + V1/V5 kernel inheritance + upstream paper results}. No primitive is forced from a deeper layer; the empirical case is cross-domain reach.

What this paper claims:

\begin{itemize}[noitemsep]
\item The canonical enumeration P01--P13 is the definitive Wave-2 specification.
\item The postulational genre --- conditional structural derivation given postulated primitives --- is the program's framing.
\item Every Wave-2 paper references this position paper for primitive enumeration.
\item The P14 placeholder is reserved openly until resolved.
\end{itemize}

What this paper does \emph{not} claim:

\begin{itemize}[noitemsep]
\item The primitives are not derived; they are postulated.
\item ED is not the only consistent substrate ontology.
\item No derivation from nothing.
\item No claim that the primitive set is minimal in some absolute sense.
\item P14 is not part of the canonical 13.
\end{itemize}

\section*{Appendix: Glossary}

\begin{itemize}[noitemsep]
\item \textbf{13 primitives.} P01--P13 as enumerated in \S4.
\item \textbf{P14 placeholder.} Reserved slot for Bilocal Strain Coupling (Paper\_030 \S2); not in canonical 13.
\item \textbf{Postulational genre.} Primitives postulated; downstream content conditional structural derivation.
\item \textbf{Cross-domain reach.} Empirical case: same 13 primitives generate mechanisms across gravity/QM/BH/soft-matter/Q-COMPUTE/entanglement.
\item \textbf{Conditional structural derivation.} Wave-2 generative paper format: given primitives + inheritance + upstream, derive structural consequence.
\item \textbf{Substrate-ontology research-program lineage.} 't Hooft cellular-automaton interpretation, Wolfram Ruliad / hypergraph rewriting, causal-set program (Sorkin et al.). ED's closest intellectual neighborhood: ontology + cross-domain exhibits, not closed-proof reconstruction.
\item \textbf{Methodological precedent (postulational genre).} Newton's three laws, Einstein's special-relativity postulates, Schr\"odinger's equation, Dirac's spinor equation --- \emph{positing axioms and arguing from downstream reach} as a research practice. ED invokes this precedent for its methodology, not for equivalent epistemic standing.
\item \textbf{Operational-reconstruction tradition (NOT ED's lineage).} Hardy 2001, CDP 2011, Masanes-M\"uller 2011, Coecke-Kissinger 2017 --- programs that derive finite-dimensional QM from operational axioms by \textbf{closed proofs}. ED is \emph{not} in this lineage; ED postulates a substrate and exhibits cross-domain reach without a closed substrate$\to$continuum reconstruction theorem.
\end{itemize}

\end{document}
